\chapter{交互作用}

再考虑一个随机化实验，它要回答的因果问题是：“抬头望天会不会让别的行人也跟着抬头？”迄今为止，我们只把兴趣限定在单一处理（抬头）在整个总体或其中某个子集上的因果效应。然而，许多因果问题实际上涉及两个或更多同时施加的处理。例如，设想除了随机分配你是否抬头，我们还随机分配你是穿戴整齐站在街上、还是一丝不挂。这时就可以问：若你穿着衣服，抬头的因果效应是什么？若你赤身裸体呢？若这两个因果效应不同，我们就说所考虑的两个处理（抬头与穿着）在导致结局上存在交互作用。

当对两个或更多处理的联合干预可行时，识别交互作用使人得以实施最有效的干预。因此，理解交互作用的概念对因果推断至关重要。本章给出两个处理之间交互作用的形式定义，既在我们已经熟悉的反事实框架内，也在充分成分原因框架内。

\section{交互作用需要联合干预}

设想在我们的心脏移植例子中，个体先被分配接受复合维生素（$E=1$）或不接受维生素（$E=0$），随后被分配接受心脏移植（$A=1$）或不接受心脏移植（$A=0$）。现在可以把所有个体分为 4 个处理组：维生素—移植（$E=1,A=1$）、维生素—不移植（$E=1,A=0$）、无维生素—移植（$E=0,A=1$）、无维生素—不移植（$E=0,A=0$）。对每个个体，现在可以设想 4 个潜在结局或反事实结局，分别对应这 4 种处理组合：$Y^{a=1,e=1}$、$Y^{a=1,e=0}$、$Y^{a=0,e=1}$ 与 $Y^{a=0,e=0}$。一般地，个体的反事实结局 $Y^{a,e}$ 是：若我们施加干预、把该个体的 $A$ 与 $E$ 分别设为 $a$ 与 $e$，本可观测到的结局。我们把对两个或更多处理的干预称为联合干预。

\mnote{只对 $A$ 施加干预所对应的反事实 $Y^{a}$，就是观测到的 $E$ 取值 $e$ 之下的联合反事实 $Y^{a,e}$，即 $Y^{a}=Y^{a,E}$。事实上，一致性正是这种递归代入的一个特例。具体而言，观测到的 $Y=Y^{A}=Y^{A,E}$，这就是我们给一致性下的定义。另见技术注记 6.2。}

现在可以在反事实框架内给出交互作用的定义。若在一次把 $E$ 设为 1 的联合干预之后，$A$ 对 $Y$ 的因果效应，不同于一次把 $E$ 设为 0 的联合干预之后 $A$ 对 $Y$ 的因果效应，则两个处理 $A$ 与 $E$ 之间存在交互作用。例如，若人人都服用维生素时移植对存活的因果效应，不同于无人服用维生素时移植对存活的因果效应，则移植 $A$ 与维生素 $E$ 之间存在交互作用。

当因果效应以风险差尺度度量时，我们说：若
\[
\begin{aligned}
&\Pr\bigl[Y^{a=1,e=1}=1\bigr]-\Pr\bigl[Y^{a=0,e=1}=1\bigr]\\
&\qquad\neq \Pr\bigl[Y^{a=1,e=0}=1\bigr]-\Pr\bigl[Y^{a=0,e=0}=1\bigr],
\end{aligned}
\]
则总体中 $A$ 与 $E$ 之间存在加法尺度上的交互作用。

例如，设人人都服用维生素时，移植 $A$ 的因果风险差 $\Pr[Y^{a=1,e=1}=1]-\Pr[Y^{a=0,e=1}=1]$ 为 0.1；无人服用维生素时移植 $A$ 的因果风险差 $\Pr[Y^{a=1,e=0}=1]-\Pr[Y^{a=0,e=0}=1]$ 为 0.2。我们说 $A$ 与 $E$ 之间存在加法尺度上的交互作用，因为风险差 $\Pr[Y^{a=1,e=1}=1]-\Pr[Y^{a=0,e=1}=1]$ 小于风险差 $\Pr[Y^{a=1,e=0}=1]-\Pr[Y^{a=0,e=0}=1]$。用简单的代数即可表明，这一不等式蕴含：人人都接受移植时维生素 $E$ 的因果风险差 $\Pr[Y^{a=1,e=1}=1]-\Pr[Y^{a=1,e=0}=1]$，也小于无人接受移植 $A$ 时维生素 $E$ 的因果风险差 $\Pr[Y^{a=0,e=1}=1]-\Pr[Y^{a=0,e=0}=1]$。也就是说，我们可以等价地把 $A$ 与 $E$ 之间加法尺度上的交互作用定义为
\[
\begin{aligned}
&\Pr\bigl[Y^{a=1,e=1}=1\bigr]-\Pr\bigl[Y^{a=1,e=0}=1\bigr]\\
&\qquad\neq \Pr\bigl[Y^{a=0,e=1}=1\bigr]-\Pr\bigl[Y^{a=0,e=0}=1\bigr].
\end{aligned}
\]
上面两个不等式表明，在交互作用的定义中，处理 $A$ 与 $E$ 地位对等。另见技术注记 5.1。

下面回顾交互作用与效应修饰的区别。如上一章所述，当 $A$ 对 $Y$ 的平均因果效应随 $V$ 的取值水平而异时，称变量 $V$ 是 $A$ 之效应的修饰因子。请注意，效应修饰的概念指的是 $A$ 的因果效应，而不是 $V$ 的因果效应。例如，在表 4.1 中性别是心脏移植效应的修饰因子，但我们从未讨论性别对死亡的效应。因此，当我们说 $V$ 修饰 $A$ 的效应时，并未把 $V$ 与 $A$ 视为地位对等的变量，因为只有 $A$ 被看作我们可以在假想中对其施加干预的变量。也就是说，效应修饰的定义涉及反事实结局 $Y^{a}$，而不是反事实结局 $Y^{a,e}$。相反，$A$ 与 $E$ 之间交互作用的定义赋予两个处理 $A$ 与 $E$ 同等地位，上面两个等价的交互作用定义正体现了这一点。交互作用的概念指的是两个处理 $A$ 与 $E$ 的联合因果效应，因而涉及联合干预之下的反事实结局 $Y^{a,e}$。

\section{识别交互作用}

前几章描述了识别处理 $A$ 对结局 $Y$ 的平均因果效应所需的条件，无论在整个总体中还是在某个子集中。三个关键的识别条件是可交换性、正值性与一致性。由于交互作用关心的是两个（或更多）处理 $A$ 与 $E$ 的联合效应，识别交互作用要求两个处理都满足可交换性、正值性与一致性。

设维生素 $E$ 由研究者随机且无条件地分配。于是正值性与一致性成立，且受维生素者 $E=1$ 与未受维生素者 $E=0$ 预期是可交换的。也就是说，若所有个体都被分配到移植 $A=1$ 与维生素 $E=1$，本可观测到的风险，等于若所有接受了 $E=1$ 的个体都被分配到移植 $A=1$ 本可观测到的风险。形式上，边缘风险 $\Pr[Y^{a=1,e=1}=1]$ 等于条件风险 $\Pr[Y^{a=1}=1\given E=1]$。因此，我们可以把 $A$ 与 $E$ 之间加法尺度上交互作用的定义改写为
\[
\begin{aligned}
&\Pr\bigl[Y^{a=1}=1\given E=1\bigr]-\Pr\bigl[Y^{a=0}=1\given E=1\bigr]\\
&\qquad\neq \Pr\bigl[Y^{a=1}=1\given E=0\bigr]-\Pr\bigl[Y^{a=0}=1\given E=0\bigr],
\end{aligned}
\]
这恰是 $E$ 对 $A$ 之效应在加法尺度上的修饰的定义。换言之，当处理 $E$ 被随机分配时，交互作用与效应修饰这两个概念重合。第 4 章所述识别 $V$ 对 $A$ 之效应修饰的方法，只需把效应修饰因子 $V$ 换成处理 $E$，即可用来识别 $A$ 与 $E$ 的交互作用。

\renewcommand{\TPname}{技术注记（Technical Point）}
\begin{techpoint}[加法尺度与乘法尺度上的交互作用]
因果风险差满足
\[
\Pr\bigl[Y^{a=1,e=1}=1\bigr]-\Pr\bigl[Y^{a=0,e=1}=1\bigr]=\Pr\bigl[Y^{a=1,e=0}=1\bigr]-\Pr\bigl[Y^{a=0,e=0}=1\bigr]
\]
这一等式可以改写为
\[
\Pr\bigl[Y^{a=1,e=1}=1\bigr]=\Bigl\{\Pr\bigl[Y^{a=1,e=0}=1\bigr]-\Pr\bigl[Y^{a=0,e=0}=1\bigr]\Bigr\}+\Pr\bigl[Y^{a=0,e=1}=1\bigr].
\]
从等式两边同时减去 $\Pr[Y^{a=0,e=0}=1]$，得到
\[
\begin{aligned}
&\Pr\bigl[Y^{a=1,e=1}=1\bigr]-\Pr\bigl[Y^{a=0,e=0}=1\bigr]\\
&\qquad=\Bigl\{\Pr\bigl[Y^{a=1,e=0}=1\bigr]-\Pr\bigl[Y^{a=0,e=0}=1\bigr]\Bigr\}\\
&\qquad\quad+\Bigl\{\Pr\bigl[Y^{a=0,e=1}=1\bigr]-\Pr\bigl[Y^{a=0,e=0}=1\bigr]\Bigr\}.
\end{aligned}
\]
这一等式是表明 $A$ 与 $E$ 在交互作用定义中地位对等的另一种紧凑方式。当上述等式成立时，我们说 $A$ 与 $E$ 之间没有加法尺度上的交互作用，并说因果风险差 $\Pr[Y^{a=1,e=1}=1]-\Pr[Y^{a=0,e=0}=1]$ 是可加的，因为它可以写成两个因果风险差之和：一个度量无 $E$ 时 $A$ 的效应，另一个度量无 $A$ 时 $E$ 的效应。反之，若
\[
\begin{aligned}
&\Pr\bigl[Y^{a=1,e=1}=1\bigr]-\Pr\bigl[Y^{a=0,e=0}=1\bigr]\\
&\qquad\neq\Bigl\{\Pr\bigl[Y^{a=1,e=0}=1\bigr]-\Pr\bigl[Y^{a=0,e=0}=1\bigr]\Bigr\}\\
&\qquad\quad+\Bigl\{\Pr\bigl[Y^{a=0,e=1}=1\bigr]-\Pr\bigl[Y^{a=0,e=0}=1\bigr]\Bigr\},
\end{aligned}
\]
则 $A$ 与 $E$ 之间存在加法尺度上的交互作用。若其中的“不等于”（$\neq$）可换成“大于”（$>$），则交互作用是超加性的；若可换成“小于”（$<$），则交互作用是次加性的。

类似地，当效应度量是因果风险比而非因果风险差时，可以定义乘法尺度上的交互作用。我们说 $A$ 与 $E$ 之间存在乘法尺度上的交互作用，若
\[
\frac{\Pr[Y^{a=1,e=1}=1]}{\Pr[Y^{a=0,e=0}=1]}
\neq
\frac{\Pr[Y^{a=1,e=0}=1]}{\Pr[Y^{a=0,e=0}=1]}
\times
\frac{\Pr[Y^{a=0,e=1}=1]}{\Pr[Y^{a=0,e=0}=1]}.
\]
若其中的“不等于”（$\neq$）可换成“大于”（$>$），则交互作用是超乘性的；若可换成“小于”（$<$），则交互作用是次乘性的。
\end{techpoint}
\renewcommand{\TPname}{技术注记}

现在设处理 $E$ 并非由研究者分配。要评估 $A$ 与 $E$ 之间是否存在交互作用，仍需要计算 4 个边缘风险 $\Pr[Y^{a,e}=1]$。在没有边缘随机化的情况下，只要通常的识别假设成立，这些风险对两个处理 $A$ 与 $E$ 都可以通过在已测协变量上做标准化或逆概率加权算出。思考这一问题的另一种等价方式是：与其把 $A$ 与 $E$ 看作各有两个可能水平（1 或 0）的两个不同处理，不如把 $AE$ 看作一个有四个可能水平（11、01、10、00）的复合处理。按这种理解，两个处理之间交互作用的识别，与前几章讨论过的单一处理因果效应的识别并无不同：同样的方法、在同样的可识别条件下均可使用。唯一的差别在于，现在所关心的处理可以取的取值列表更长，因而反事实结局的数目更多。

有时人们愿意对处理 $A$ 假定（条件）可交换性，却不愿对处理 $E$ 作此假定，例如在随机化实验中估计 $A$ 在由 $E$ 定义的各亚组中的因果效应时。此时一般无法评估 $A$ 与 $E$ 之间是否存在交互作用，但仍可评估 $E$ 的效应修饰是否存在。之所以如此，是因为计算 $A$ 在 $E$ 所定义的每一层内的效应，并不需要任何涉及 $E$ 的识别假设。上一章我们用记号 $V$（而不是 $E$）表示那些我们不愿就其可交换性、正值性与一致性作出假定的变量。例如，我们曾断定移植 $A$ 的效应被国籍 $V$ 修饰，但从未要求关于 $V$ 之效应的任何识别假设，因为我们并不想用数据去计算 $V$ 对 $Y$ 的因果效应。在第 4.2 节中，我们基于实质理由论证 $V$ 是一个替代效应修饰因子；也就是说，$V$ 不作用于结局，因而不会与 $A$ 发生交互作用——无作用，即无交互。但 $V$ 是 $A$ 对 $Y$ 之效应的修饰因子，因为 $V$ 与某个真正作用于 $Y$ 且与 $A$ 交互的未识别变量相关（例如是它的替代指标）。因此，$A$ 的效应可以被另一个变量修饰，而 $A$ 与该变量之间并无交互作用。

\mnote{$A$ 与 $E$ 之间存在交互作用而 $E$ 并不修饰 $A$ 的效应，在逻辑上也是可能的，尽管大概很罕见，因为它要求 $A$ 具有双重效应并且发生精确的抵消（VanderWeele 2009b）。}

以上我们论证了：识别两个处理 $A$ 与 $E$ 之间交互作用的一个充分条件是，联合处理 $(A,E)$ 的四个可能取值 $(0,0)$、$(0,1)$、$(1,0)$、$(1,1)$ 全都满足可交换性、正值性与一致性。于是可以用标准化或逆概率加权估计两个处理的联合效应，进而评估它们之间的交互作用。在第三部分我们将说明，当两个处理发生在不同时点时，这一条件并非必要。就第一部分其余各章（本章除外）与第二部分的大部分而言，我们仍将聚焦于单一处理 $A$ 的因果效应。

第 1 章描述了确定性反事实结局与非确定性反事实结局。至此为简单起见，我们用的都是确定性反事实。然而，我们讨论过的关于总体因果效应与交互作用的结果，没有一项要求反事实结局是确定性的。相反，本章下一节只适用于反事实为确定性的情形。此外，我们还假定处理与结局都是二值的。

\section{反事实反应类型与交互作用}

可以按确定性的反事实反应对个体分类。例如，在表 4.1（与表 1.1 相同）中，人有四种类型：“注定型”，不论接受何种处理都会发生结局（Artemis、Athena、Persephone、Ares）；“免疫型”，不论接受何种处理都不会发生结局（Demeter、Hestia、Hera、Hades）；“受益型”，只有未受治疗时才会发生结局（Hebe、Kronos、Poseidon、Apollo、Hermes、Dyonisus）；以及“受害型”，只有受治疗时才会发生结局（Rheia、Leto、Aphrodite、Zeus、Hephaestus、Polyphemus）。反事实反应的每一种组合，常被称为一种反应模式或反应类型。表 5.1 给出四种可能的反应类型。

\begin{table}[htbp]
\centering
\caption{四种可能的反应类型}
\label{tab:5.1}
\small
\begin{tabular}{lcc}
\toprule
类型 & $Y^{a=0}$ & $Y^{a=1}$\\
\midrule
注定型 & 1 & 1\\
受益型 & 1 & 0\\
受害型 & 0 & 1\\
免疫型 & 0 & 0\\
\bottomrule
\end{tabular}
\end{table}

考虑两个二值处理 $A$ 与 $E$ 时，共有 16 种可能的反应类型，因为每个个体有 4 个反事实结局，分别对应处理 $A$ 与 $E$ 的 4 种可能联合干预 $(1,1)$、$(0,1)$、$(1,0)$、$(0,0)$。表 5.2 给出两个处理的 16 种反应类型。本节在两个二值处理 $A$ 与 $E$ 及二值结局 $Y$ 的情形下，探讨反应类型与交互作用是否存在之间的关系。

\begin{table}[htbp]
\centering
\caption{两个处理下的 16 种反应类型}
\label{tab:5.2}
\small
\renewcommand{\arraystretch}{0.98}
\begin{tabular}{ccccc}
\toprule
\multirow{2}{*}{类型} & \multicolumn{4}{c}{各 $a,e$ 取值下的 $Y^{a,e}$}\\
\cmidrule(lr){2-5}
& $1,1$ & $0,1$ & $1,0$ & $0,0$\\
\midrule
1  & 1 & 1 & 1 & 1\\
2  & 1 & 1 & 1 & 0\\
3  & 1 & 1 & 0 & 1\\
4  & 1 & 1 & 0 & 0\\
5  & 1 & 0 & 1 & 1\\
6  & 1 & 0 & 1 & 0\\
7  & 1 & 0 & 0 & 1\\
8  & 1 & 0 & 0 & 0\\
9  & 0 & 1 & 1 & 1\\
10 & 0 & 1 & 1 & 0\\
11 & 0 & 1 & 0 & 1\\
12 & 0 & 1 & 0 & 0\\
13 & 0 & 0 & 1 & 1\\
14 & 0 & 0 & 1 & 0\\
15 & 0 & 0 & 0 & 1\\
16 & 0 & 0 & 0 & 0\\
\bottomrule
\end{tabular}
\end{table}

表 5.2 中第一种类型的反事实结局 $Y^{a=1,e=1}$ 等于 1，这意味着这一类型的个体若同时接受移植与维生素就会死亡。其余三个反事实结局也等于 1，即 $Y^{a=1,e=1}=Y^{a=0,e=1}=Y^{a=1,e=0}=Y^{a=0,e=0}=1$，这意味着这一类型的个体若接受（不移植，维生素）、（移植，无维生素）或（不移植，无维生素），同样会死亡。换言之，处理 $A$ 与处理 $E$ 对这类个体的结局都没有任何效应。无论被分配到何种联合处理，他都会死。现在考虑类型 16。所有反事实结局都是 0，即 $Y^{a=1,e=1}=Y^{a=0,e=1}=Y^{a=1,e=0}=Y^{a=0,e=0}=0$。同样，处理 $A$ 与处理 $E$ 对这一类型个体的结局都没有任何效应。无论被分配到何种联合处理，她都会存活。若总体中所有个体都属于类型 1 与 16，我们就会说 $A$ 与 $E$ 对 $Y$ 都没有任何因果效应；对联合处理 $(A,E)$ 而言，严格因果零假设成立。

下面把注意力集中在类型 4、6、11 与 13 上。类型 4 的个体只有在接受维生素时才会死亡，无论是否接受移植，即 $Y^{a=1,e=1}=Y^{a=0,e=1}=1$ 且 $Y^{a=1,e=0}=Y^{a=0,e=0}=0$。类型 13 的个体只有在不接受维生素时才会死亡，无论是否接受移植，即 $Y^{a=1,e=1}=Y^{a=0,e=1}=0$ 且 $Y^{a=1,e=0}=Y^{a=0,e=0}=1$。类型 6 的个体只有在接受移植时才会死亡，无论是否接受维生素，即 $Y^{a=1,e=1}=Y^{a=1,e=0}=1$ 且 $Y^{a=0,e=1}=Y^{a=0,e=0}=0$。类型 11 的个体只有在不接受移植时才会死亡，无论是否接受维生素，即 $Y^{a=1,e=1}=Y^{a=1,e=0}=0$ 且 $Y^{a=0,e=1}=Y^{a=0,e=0}=1$。

\mnote{Miettinen (1982) 描述了两种二值处理与二值结局下的 16 种可能反应类型。}

在表 5.2 的 16 种可能反应类型中，我们找出了 6 种（编号 1、4、6、11、13、16）具有共同特征的类型：对属于这些反应类型之一的个体，处理 $A$ 对结局 $Y$ 的因果效应与处理 $E$ 的取值无关，处理 $E$ 对结局 $Y$ 的因果效应也与处理 $A$ 的取值无关。在一个每个个体都属于这 6 种反应类型之一的总体中，存在处理 $E$ 时处理 $A$ 的因果效应（以因果风险差 $\Pr[Y^{a=1,e=1}=1]-\Pr[Y^{a=0,e=1}=1]$ 度量），将等于不存在处理 $E$ 时处理 $A$ 的因果效应（以因果风险差 $\Pr[Y^{a=1,e=0}=1]-\Pr[Y^{a=0,e=0}=1]$ 度量）。也就是说，若总体中所有个体都属于反应类型 1、4、6、11、13 与 16，则 $A$ 与 $E$ 之间不存在加法尺度上的交互作用。

\mnote{Greenland 与 Poole (1988) 指出，Miettinen 的反应类型对 $A$ 与 $E$ 的重新编码（即把标签“0”与“1”互换）并不保持不变。他们把表 5.2 的 16 种反应类型划归为三个对重新编码保持不变的等价类。}

$A$ 与 $E$ 之间存在加法交互作用，意味着总体中有些个体在 $A=a$ 之下的两个反事实结局的取值，若不借助 $E$ 的取值便无法确定，反之亦然。也就是说，总体中必定存在属于以下三类中至少一类的个体：

\begin{enumerate}
\item 只在四种处理组合之一的情况下发生结局的个体（表 5.2 中的类型 8、12、14 与 15）；
\item 在两种处理组合下发生结局、且每种处理的效应在另一处理的每一水平下恰好相反的个体（类型 7 与 10）；
\item 在四种处理组合中的三种下发生结局的个体（类型 2、3、5 与 9）。
\end{enumerate}

另一方面，$A$ 与 $E$ 之间不存在加法交互作用，意味着要么总体中没有任何个体属于上述三类，要么对加性的偏离大小相等、符号相反，恰好完全抵消。例如，若类型 7 与类型 10 的个体比例相等，或类型 8 与类型 12 的个体比例相等，就会发生这种抵消。

\mnote{关于交互作用类型存在时仍出现可加性抵消的更多讨论，见 Greenland、Lash 与 Rothman (2008)。}

“交互作用”一词的含义，通过按反事实反应类型对个体分类而变得清晰（另见精要 5.1）。下面我们引入一个工具，用来表示两个处理交互作用所涉及的因果机制。

\section{充分原因}

交互作用的含义，通过按反事实反应类型对个体分类而变得清晰。现在我们引入一个工具，用来表示两个处理交互作用所涉及的因果机制。再考虑只有一个处理 $A$ 的心脏移植例子。如上一节所述，有些个体受治疗时会死，有些未受治疗时会死，有些无论怎样都会死，还有些无论怎样都不会死。反应类型如此多样，说明处理 $A$ 并非决定结局 $Y$ 是否发生的唯一变量。

取那些实际受治疗的个体。其中只有一部分人死亡，这意味着单靠治疗不足以保证结局发生。作为一个过度简化的例子，假设心脏移植 $A=1$ 只在对接麻药过敏的个体中导致死亡。我们把与 $A=1$ 合在一起足以不可避免地产生结局的最小背景因素集合记作 $U_1$。治疗（$A=1$）与对接麻药过敏（$U_1=1$）同时存在，构成结局 $Y$ 的一个\emph{最小充分原因}。

再取那些未受治疗的个体。同样只有一部分人死亡，这意味着单靠不治疗也不足以导致结局。作为一个过度简化的例子，假设不施行心脏移植 $A=0$ 只在射血分数低于 20\% 的个体中导致死亡。我们把与 $A=0$ 合在一起足以产生结局的最小背景因素集合记作 $U_2$。不治疗（$A=0$）与低射血分数（$U_2=1$）同时存在，是结局 $Y$ 的另一个充分原因。

\renewcommand{\FPname}{精要（Fine Point）}
\begin{finepoint}[关于反事实类型与交互作用的补充]
按反事实反应类型对个体分类，使我们可以更方便地考虑特定形式的交互作用。例如，我们可能想知道是否存在这样的个体：同时接受两种处理 $E=1$ 与 $A=1$ 时会发生结局，而只接受其中之一时不会。也就是说，总体中是否存在反事实反应为 $Y^{a=1,e=1}=1$ 且 $Y^{a=0,e=1}=Y^{a=1,e=0}=0$ 的个体（类型 7 与 8）。VanderWeele 与 Robins (2007a, 2008) 针对 2 个与 3 个处理建立了充分原因交互作用的理论，并导出了此处所述的协同作用识别条件。下述不等式是这些个体存在的一个充分条件：
\[
\Pr\bigl[Y^{a=1,e=1}=1\bigr]-\Bigl(\Pr\bigl[Y^{a=0,e=1}=1\bigr]+\Pr\bigl[Y^{a=1,e=0}=1\bigr]\Bigr)>0
\]
或等价地，
\[
\Pr\bigl[Y^{a=1,e=1}=1\bigr]-\Pr\bigl[Y^{a=0,e=1}=1\bigr]>\Pr\bigl[Y^{a=1,e=0}=1\bigr].
\]
也就是说，在 $A$ 与 $E$ 被随机分配的实验中，人们可以算出上述不等式中的三个反事实风险，从而凭经验检验类型 7 与 8 的个体是否存在。

由于上述不等式只是充分而非必要条件，即使类型 7 与 8 存在，它也可能不成立。事实上，这个充分条件如此之强，以致可能漏掉这些类型存在的大多数情形。若已知或愿意假定接受处理 $A$ 与 $E$ 不能阻止任何个体发生结局，即效应是单调的（见技术注记 5.2），则可以使用一个较弱的协同作用充分条件。此时，不等式
\[
\begin{aligned}
&\Pr\bigl[Y^{a=1,e=1}=1\bigr]-\Pr\bigl[Y^{a=0,e=1}=1\bigr]\\
&\qquad>\Pr\bigl[Y^{a=1,e=0}=1\bigr]-\Pr\bigl[Y^{a=0,e=0}=1\bigr]
\end{aligned}
\]
是类型 7 与 8 存在的一个充分条件。换言之，当 $A$ 与 $E$ 的效应单调时，超加性交互作用的存在蕴含类型 8 的存在（单调性排除类型 7）。这一单调效应下协同作用的充分条件，最初由 Greenland 与 Rothman 在他们著作的前一版中报告，现载于 Greenland、Lash 与 Rothman (2008)。

在遗传学研究中，有时需要确定是否存在类型 8 的个体，这种交互作用形式称为组合性上位效应。VanderWeele (2010a) 综述了组合性上位效应的经验检验。
\end{finepoint}
\renewcommand{\FPname}{精要}

最后，假设有些个体既无 $U_1$ 也无 $U_2$，且无论受治疗与否都会发生结局。这些“注定型”个体的存在意味着另有一些背景因素本身即足以导致结局。作为一个过度简化的例子，假设研究开始时已患胰腺癌的个体都会死亡。我们把不论治疗状态如何都足以产生结局的最小背景因素集合记作 $U_0$。患胰腺癌（$U_0=1$）是结局 $Y$ 的又一个充分原因。

\mnote{按背景因素的定义，这些二值变量 $U$ 不能被干预，也不会受处理 $A$ 的影响。}

我们描述了结局的 3 个充分原因：存在 $U_1$ 时的治疗 $A=1$、存在 $U_2$ 时的不治疗 $A=0$，以及不论治疗状态如何的 $U_0$ 存在。每个充分原因都有一个或多个\emph{成分}。例如，第一个充分原因的成分为 $A=1$ 与 $U_1=1$。图 5.1 用一个圆表示每个充分原因，用圆的分块表示其成分。人们常用“充分成分原因”一词来指称充分原因及其成分。

\begin{figure}[htbp]
\centering
\begin{tikzpicture}[causal,scale=1]
  \begin{scope}
    \draw[thick] (0,0) circle (0.8);
    \draw[thick] (-0.8,0)--(0.8,0);
    \node[font=\small] at (0,0.4) {$A=1$};
    \node[font=\small] at (0,-0.4) {$U_1=1$};
  \end{scope}
  \begin{scope}[xshift=2.8cm]
    \draw[thick] (0,0) circle (0.8);
    \draw[thick] (-0.8,0)--(0.8,0);
    \node[font=\small] at (0,0.4) {$A=0$};
    \node[font=\small] at (0,-0.4) {$U_2=1$};
  \end{scope}
  \begin{scope}[xshift=5.6cm]
    \draw[thick] (0,0) circle (0.8);
    \node[font=\small] at (0,0) {$U_0=1$};
  \end{scope}
\end{tikzpicture}
\caption{结局 $Y$ 的三个充分原因及其成分}
\label{fig:5.1}
\end{figure}

充分成分原因的图示有助于直观理解效应修饰的一个关键推论：如第 4 章所论，处理 $A$ 的因果效应大小取决于效应修饰因子的分布。设想两个假想情形。第一个情形中，总体里只有 1\% 的个体有 $U_1=1$（即对接麻药过敏）。第二个情形中，总体里有 10\% 的个体有 $U_1=1$。两个总体中 $U_2$ 与 $U_0$ 的分布相同。现在，分别在每个总体中开展心脏移植 $A$ 的随机化实验，其中一半人被分配到治疗 $A=1$。心脏移植 $A$ 对死亡的平均因果效应在第二个总体中更大，因为那里有更多个体在受治疗时易感于发生结局。3 个充分原因之一（$A=1$ 加 $U_1=1$）在第二个总体中的常见程度是第一个总体的 10 倍，而另外两个充分原因在两个总体中同样常见。

充分成分原因的图示还有助于直观理解交互作用的另一种概念，下一节将加以描述。首先需要描述两个处理 $A$ 与 $E$ 的充分原因。考虑我们的维生素与心脏移植例子。我们已经描述了死亡的 3 个充分原因：其一与 $A$（或 $E$）是否存在无关，其二是不论维生素 $E$ 如何都存在移植 $A$，其三是不论维生素 $E$ 如何都不存在移植 $A$。在两个处理的情形下，还需再加 2 种死亡途径：不论移植 $A$ 如何都存在维生素 $E$，以及不论移植 $A$ 如何都不存在维生素 $E$。此外还需再加四个充分原因，以容纳那些只在处理 $A$ 与 $E$ 的某些取值组合下才会死亡的个体。于是，取决于存在哪些背景因素，共有 9 种可能的死亡途径：

\begin{enumerate}
\item 因处理 $A$ 而死（处理 $E$ 无关）
\item 因缺少处理 $A$ 而死（处理 $E$ 无关）
\item 因处理 $E$ 而死（处理 $A$ 无关）
\item 因缺少处理 $E$ 而死（处理 $A$ 无关）
\item 因同时接受处理 $A$ 与 $E$ 而死
\item 因接受处理 $A$ 且缺少 $E$ 而死
\item 因接受处理 $E$ 且缺少 $A$ 而死
\item 因同时缺少 $A$ 与 $E$ 而死
\item 因其他机制而死（处理 $A$ 与 $E$ 均无关）
\end{enumerate}

换言之，共有 9 个可能的充分原因，其处理成分分别为：仅 $A=1$、仅 $A=0$、仅 $E=1$、仅 $E=0$、$A=1$ 与 $E=1$、$A=1$ 与 $E=0$、$A=0$ 与 $E=1$、$A=0$ 与 $E=0$，以及 $A$ 与 $E$ 均无关。每个这样的充分原因都包含一组来自 $U_1,\dots,U_8$ 与 $U_0$ 的背景因素。图 5.2 表示两个处理 $A$ 与 $E$ 的 9 个充分成分原因。

\mnote{Greenland 与 Poole (1988) 首次枚举了这 9 个充分原因。}

\newcommand{\sccOne}[1]{%
  \draw[thick] (0,0) circle (0.7);
  \node[font=\footnotesize] at (0,0) {$#1$};}
\newcommand{\sccTwo}[2]{%
  \draw[thick] (0,0) circle (0.7);
  \draw[thick] (0,0.7)--(0,-0.7);
  \node[font=\footnotesize,anchor=east] at (-0.75,0.16) {$#1$};
  \node[font=\footnotesize,anchor=west] at (0.75,0.16) {$#2$};}
\newcommand{\sccThree}[3]{%
  \draw[thick] (0,0) circle (0.7);
  \draw[thick] (0,0)--(0,0.7);
  \draw[thick] (0,0)--(-0.606,-0.35);
  \draw[thick] (0,0)--(0.606,-0.35);
  \node[font=\footnotesize,anchor=east] at (-0.73,0.2) {$#1$};
  \node[font=\footnotesize,anchor=west] at (0.73,0.2) {$#2$};
  \node[font=\footnotesize,anchor=north] at (0,-0.7) {$#3$};}

\begin{figure}[htbp]
\centering
\begin{tikzpicture}[causal,scale=1]
  \begin{scope}
    \sccOne{U_0=1}
  \end{scope}
  \begin{scope}[xshift=4.4cm]
    \sccTwo{A=1}{U_1=1}
  \end{scope}
  \begin{scope}[xshift=9.2cm]
    \sccTwo{A=0}{U_2=1}
  \end{scope}
  \begin{scope}[yshift=-2.4cm]
    \sccTwo{U_3=1}{E=1}
  \end{scope}
  \begin{scope}[xshift=4.4cm,yshift=-2.4cm]
    \sccTwo{U_4=1}{E=0}
  \end{scope}
  \begin{scope}[xshift=9.2cm,yshift=-2.4cm]
    \sccThree{A=1}{E=1}{U_5=1}
  \end{scope}
  \begin{scope}[yshift=-4.8cm]
    \sccThree{A=1}{E=0}{U_6=1}
  \end{scope}
  \begin{scope}[xshift=4.4cm,yshift=-4.8cm]
    \sccThree{A=0}{E=1}{U_7=1}
  \end{scope}
  \begin{scope}[xshift=9.2cm,yshift=-4.8cm]
    \sccThree{A=0}{E=0}{U_8=1}
  \end{scope}
\end{tikzpicture}
\caption{两个处理 $A$ 与 $E$ 的 9 个充分成分原因}
\label{fig:5.2}
\end{figure}

\mnote{充分成分原因的这种图示，常被称为“因果饼图”。}

对二值结局与两个处理而言，这 9 个充分成分原因并非在任何场合都存在。例如，若接受维生素 $E=1$ 不会导致任何个体死亡，无论其处理 $A$ 如何，则带有成分 $E=1$ 的那 3 个充分原因就不会存在。这 3 个充分原因若存在，就意味着有些个体（例如 $U_3=1$ 者）会因接受维生素（$E=1$）而死亡，也就是说，不给他们维生素（$E=0$）便可避免其死亡。另见技术注记 5.3。

\section{充分原因交互作用}

人们口语中所说的“处理 $A$ 与 $E$ 之间的交互作用”，唤起的是这样一种意象：存在某种因果机制，使两个处理共同起作用（即“相互作用”）以产生某个结局。有趣的是，反事实框架内对交互作用的定义，既不要求了解这些机制，甚至不要求两个处理真的共同起作用（见精要 5.3）。在维生素 $E$ 与心脏移植 $A$ 的例子中我们说，若人人接受 $E$ 时 $A$ 的因果效应不同于无人接受 $E$ 时 $A$ 的因果效应，则处理 $A$ 与 $E$ 之间存在交互作用。也就是说，交互作用由反事实量之间的对比来定义，因此可以通过开展一个理想随机化实验来识别——只要可交换性、正值性与一致性对两个处理 $A$ 与 $E$ 都成立。无需考虑交互作用存在背后的因果机制（物理的、化学的、生物的、社会的……）。

本节描述交互作用的第二种概念，它或许使我们离处理 $A$ 与 $E$ 产生结局的因果机制更近一步。这第二种交互作用概念不基于反事实对比，而基于充分成分原因，因此我们称之为充分成分原因框架内的交互作用，或简称为\emph{充分原因交互作用}。

若 $A$ 与 $E$ 出现在同一个充分原因中，则总体中存在 $A$ 与 $E$ 之间的充分原因交互作用。例如，假设带有背景因素 $U_5=1$ 的个体在同时接受维生素（$E=1$）与心脏移植（$A=1$）时会发生结局，而只接受两者之一时不会。那么，只要存在一个 $U_5=1$ 的个体，$A$ 与 $E$ 之间就存在充分原因交互作用。由此可知，若存在一个反事实反应为 $Y^{a=1,e=1}=1$ 且 $Y^{a=0,e=1}=Y^{a=1,e=0}=0$ 的个体，则 $A$ 与 $E$ 之间存在充分原因交互作用。

充分原因交互作用可以是协同的，也可以是拮抗的。当 $A=1$ 与 $E=1$ 出现在同一个充分原因中时，处理 $A$ 与处理 $E$ 之间存在\emph{协同作用}；当 $A=1$ 与 $E=0$（或 $A=0$ 与 $E=1$）出现在同一个充分原因中时，二者之间存在\emph{拮抗作用}。也可以把处理 $A$ 与处理 $E$ 之间的拮抗作用理解为处理 $A$ 与“不施加处理 $E$”之间的协同作用（或“不施加处理 $A$”与处理 $E$ 之间的协同作用）。

与反事实的交互作用定义不同，充分原因交互作用明确指涉涉及处理 $A$ 与 $E$ 的因果机制。人们或许因此以为，识别充分原因交互作用是否存在，需要关于这些因果机制的详尽知识。事实并非总是如此：有时即使我们对充分原因及其成分一无所知，也能断定存在充分原因交互作用。

\renewcommand{\FPname}{精要（Fine Point）}
\begin{finepoint}[从反事实到充分成分原因，以及反过来]
反事实反应类型与充分成分原因之间存在对应关系。在二值处理与二值结局的情形下，设某个个体不具备任何背景因素 $U_0$、$U_1$、$U_2$。她将属于“免疫型”反应类型，因为无论是否受治疗，她都缺少完成所有充分原因所需的成分。下表给出一个处理 $A$ 的情形下反应类型与充分成分原因之间的对应关系。

\begin{center}
\footnotesize
\begin{tabular}{lccc}
\toprule
类型 & $Y^{a=0}$ & $Y^{a=1}$ & 成分原因\\
\midrule
注定型 & 1 & 1 & $U_0=1$ 或 $\{U_1=1$ 且 $U_2=1\}$\\
受益型 & 1 & 0 & $U_0=0$ 且 $U_1=0$ 且 $U_2=1$\\
受害型 & 0 & 1 & $U_0=0$ 且 $U_1=1$ 且 $U_2=0$\\
免疫型 & 0 & 0 & $U_0=0$ 且 $U_1=0$ 且 $U_2=0$\\
\bottomrule
\end{tabular}
\end{center}

成分原因的某一特定组合对应且仅对应一种反事实类型。然而，某一特定反应类型可以对应若干种成分原因组合。例如，“注定型”个体可以具有任何包含 $U_0=1$ 的成分原因组合，无论 $U_1$ 与 $U_2$ 取何值；也可以具有任何包含 $\{U_1=1$ 且 $U_2=1\}$ 的组合。

充分成分原因还可以用来给可交换性 $Y^{a}\indep A$ 提供一种机制性描述。对二值处理与二值结局，可交换性意味着：在受治疗者 $A=1$ 与未受治疗者 $A=0$ 中，若受治疗本会发生结局的个体比例、以及若未受治疗本会发生结局的个体比例，都相同。即 $\Pr[Y^{a=1}=1\given A=1]=\Pr[Y^{a=1}=1\given A=0]$ 且 $\Pr[Y^{a=0}=1\given A=1]=\Pr[Y^{a=0}=1\given A=0]$。

现在，若受治疗本会发生结局的个体，就是“注定型”与“受害型”，即 $U_0=1$ 或 $U_1=1$ 者；若未受治疗本会发生结局的个体，就是“注定型”与“受益型”，即 $U_0=1$ 或 $U_2=1$ 者。因此，若“注定型”加“受害型”的比例与“注定型”加“受益型”的比例在受治疗者与未受治疗者中分别相等，就存在可交换性。也就是说，二值处理与二值结局下的可交换性可以用充分成分原因表述为 $\Pr[U_0=1\ \text{或}\ U_1=1\given A=1]=\Pr[U_0=1\ \text{或}\ U_1=1\given A=0]$，且 $\Pr[U_0=1\ \text{或}\ U_2=1\given A=1]=\Pr[U_0=1\ \text{或}\ U_2=1\given A=0]$。

更多细节见 Greenland 与 Brumback (2002)、Flanders (2006) 以及 VanderWeele 与 Hernán (2006)。上述部分结果由 VanderWeele 与 Robins (2008) 推广到两个或更多二值处理的情形。
\end{finepoint}
\renewcommand{\FPname}{精要}

\renewcommand{\FPname}{精要（Fine Point）}
\begin{finepoint}[生物学交互作用]
在流行病学的讨论中，充分原因交互作用通常被称为生物学交互作用（Rothman et al, 1980）。这一术语选择似乎暗示：在生物医学应用中，存在两个处理 $A$ 与 $E$ 彼此作用从而导致结局的生物学机制。然而，正如 VanderWeele 与 Robins (2007a) 提出的下例所示，情况未必如此。

设 $A$ 与 $E$ 是某个产生必需蛋白质的基因的两个等位基因。两个等位基因都带有害突变的个体（$A=1$ 且 $E=1$）将缺乏该必需蛋白质，出生后一周内死亡；而两个等位基因都无突变（即 $A=0$ 且 $E=0$）、或只有一个等位基因有突变（即 $A=0$ 且 $E=1$、$A=1$ 且 $E=0$）的个体，该蛋白质水平正常，能够存活。我们会说等位基因 $A$ 与 $E$ 之间存在协同作用，因为存在一个包含 $A=1$ 与 $E=1$ 的死亡充分成分原因。也就是说，两个等位基因共同起作用以产生结局。然而，也可以争辩说它们在物理上并不相互作用，因而在任何生物学意义上都不存在交互作用。
\end{finepoint}
\renewcommand{\FPname}{精要}

具体而言，若精要 5.1 中的不等式成立，则 $A$ 与 $E$ 之间存在协同作用。也就是说，人们可以凭经验检验协同作用是否存在，而完全不必考虑 $A$ 与 $E$ 共同作用以产生结局的因果机制。

\mnote{Rothman (1976) 在充分成分原因框架内描述了协同作用与拮抗作用的概念。}

这一结果并不太令人意外，因为反事实反应类型与充分原因之间存在对应关系（见精要 5.2），而且上述不等式只是充分而非必要条件，即协同作用存在时该不等式也可能不成立。

\renewcommand{\TPname}{技术注记（Technical Point）}
\begin{techpoint}[因果效应的单调性]
考虑一个二值处理 $A$ 与二值结局 $Y$ 的情形。只有在“受益型”个体中，反事实结局 $Y^{a=0}$ 的取值才大于 $Y^{a=1}$ 的取值。对其余 3 种类型，$Y^{a=1}\ge Y^{a=0}$；等价地说，个体的反事实结局随 $a$ 单调递增（即不减）。因此，当治疗不能阻止任何个体发生结局时（即不存在“受益型”个体时），所有个体的反事实反应类型都随 $a$ 单调递增。这时我们就简单地说，$A$ 对 $Y$ 的因果效应是单调的。

单调性的概念可以推广到两个处理 $A$ 与 $E$。若每个个体的反事实结局 $Y^{a,e}$ 对 $a$ 与 $e$ 都单调递增，则称 $A$ 与 $E$ 对 $Y$ 的因果效应是单调的。也就是说，不存在反应类型为 $\bigl(Y^{a=1,e=1}=0,\ Y^{a=0,e=1}=1\bigr)$、$\bigl(Y^{a=1,e=1}=0,\ Y^{a=1,e=0}=1\bigr)$、$\bigl(Y^{a=1,e=0}=0,\ Y^{a=0,e=0}=1\bigr)$ 与 $\bigl(Y^{a=0,e=1}=0,\ Y^{a=0,e=0}=1\bigr)$ 的个体。
\end{techpoint}
\renewcommand{\TPname}{技术注记}

\section{反事实还是充分成分原因？}

充分成分原因框架与反事实（潜在结局）框架回答的是不同的问题。\mnote{充分成分原因框架由哲学家 Mackie (1965) 创立。他引入了 $Y$ 的 INUS 条件概念：一个条件之\emph{不充分}但\emph{必要}的组成部分，而该条件本身对 $Y$ 而言\emph{不必要}却\emph{充分}。}充分成分原因模型考虑的是这样一些行动、事件或自然状态之集合：它们合在一起不可避免地导致所考虑的结局。该模型说明的是某一特定效应的原因。它针对的问题是：“给定某一特定效应，有哪些不同事件可能是它的原因？”潜在结局或反事实模型则聚焦于某个特定原因或干预，说明该原因的各种效应。与充分成分原因框架不同，潜在结局框架针对的问题是：“若对某个特定因素施加干预、从而把它设为其实际取值之外的另一个水平，会发生什么？”与充分成分原因框架不同，反事实框架不要求详尽了解该因素影响结局的机制。

\mnote{因果推断的反事实框架，休谟（Hume, 1748）已有所暗示。}

反事实进路回答“发生了什么”，充分成分原因进路回答“它是如何发生的”。就本书的内容——估计假想干预之平均因果效应的条件与方法——而言，反事实框架是自然的框架。充分成分原因框架则有助于思考产生某一特定结局时起作用的因果机制。充分成分原因在因果推断的教学中自有其正当位置，因为它们有助于理解若干关键概念，例如因果效应的大小对背景因素（效应修饰因子）分布的依赖，以及效应修饰、交互作用与协同作用之间的关系。

尽管充分成分原因框架在教学上很有用，它与实际数据分析的相关性仍待确定。在其经典形式中，充分成分原因框架是确定性的，其结论依赖于结局的编码方式，并且按定义只限于二值处理与二值结局（或可重新编码为二值变量的变量）。这一限制实际上排除了对任何连续因素的考虑，并把该框架的适用范围限制在只有少数几个二值因素的场合。晚近把充分成分原因框架推广到随机情形以及分类与有序处理的工作，或许会使这一进路更多地应用于现实数据分析。最后，即便考虑到这些推广，我们也很少拥有研究该框架所作精细区分所需的大量数据。

要更一般地估计因果效应，反事实框架很可能仍是最常被采用的框架。一些看似另类的框架——因果图、决策论——本质上等价于反事实框架，下一章将加以描述。

\renewcommand{\FPname}{精要（Fine Point）}
\begin{finepoint}[关于归因分数的补充]
精要 3.5 把处理 $A$ 的超额分数定义为特定总体中归因于处理 $A$ 的病例所占比例，并描述了一个例子，其中 $A$ 的超额分数为 75\%。也就是说，若人人都接受处理 $a=0$ 而非其实际接受的处理 $A$，75\% 的病例本不会发生。现在考虑第二种处理 $E$。设 $E$ 的超额分数为 50\%。这是否意味着对 $A$ 与 $E$ 的联合干预可以预防 125\%（75\% 加 50\%）的病例？当然不是。

显然，对单一处理（$A$ 或 $E$）而言，超额分数不会超过 100\%。同样显然的是，对 $A$ 与 $E$ 的任何联合干预，超额分数也不会超过 100\%。也就是说，若允许我们以任何希望的方式（修改 $A$、修改 $E$，或两者都修改）对一个总体施加干预，我们所能预防的疾病比例绝不会超过 100\%。换言之，归因于缺少某种干预（无论单一干预还是联合干预）的病例不会超过 100\%。那么，为什么两个单一处理的超额分数之和会大于 100\%？充分成分原因框架有助于回答这一问题。

举例来说，设 Zeus 具有背景因素 $U_5=1$（且没有其他背景因素），并同时接受了 $A=1$ 与 $E=1$。若扣留处理 $A$ 或处理 $E$，Zeus 都不会成为病例。因此，Zeus 被计入“通过设置 $a=0$ 的干预而得以预防的病例”，即 Zeus 属于归因于 $A$ 的那 75\% 的病例。但 Zeus 也被计入“通过设置 $e=0$ 的干预而得以预防的病例”，即 Zeus 属于归因于 $E$ 的那 50\% 的病例。难怪 $A$ 与 $E$ 的超额分数之和超过 100\%：有些像 Zeus 这样的个体被数了两次！

充分成分原因框架表明，当 $A$ 与 $E$ 都可能是同一个充分原因的组成成分时，分别谈论归因于 $A$ 与归因于 $E$ 的疾病比例没有多大意义。例如，关于疾病比例应归因于基因还是环境的讨论具有误导性。考虑苯丙酮尿症所致的精神发育迟缓——这种病症出现在食用某些食物的遗传易感个体中。对那些食物的超额分数为 100\%，因为只要从饮食中去除这些食物，所有病例都可预防。对基因的超额分数也是 100\%，因为只要能替换易感基因，所有病例都可预防。于是，精神发育迟缓的原因既可以看作 100\% 遗传的，也可以看作 100\% 环境的。进一步讨论见 Rothman、Greenland 与 Lash (2008)。
\end{finepoint}
\renewcommand{\FPname}{精要}

\mnote{VanderWeele (2010b) 提供了推广到三水平处理的方法。VanderWeele 与 Robins (2012) 探讨了随机反事实与随机充分原因之间的关系。}

\renewcommand{\TPname}{技术注记（Technical Point）}
\begin{techpoint}[因果效应的单调性与充分原因]
当处理 $A$ 与 $E$ 的效应单调时，某些充分原因必定不存在。例如，设吸烟（$A=1$）从不预防心脏病，缺乏体力活动（$E=1$）也从不预防心脏病。那么，任何包含 $A=0$ 或 $E=0$ 的充分原因都不可能存在。之所以如此，是因为若存在包含成分 $A=0$ 的充分原因，则有些个体（例如 $U_2=1$ 者）在未暴露（$A=0$）时会发生结局；等价地说，对这些个体施以治疗（$A=1$）即可预防该结局。同样的道理适用于 $E=0$。当 $A$ 与 $E$ 的效应单调时不可能存在的那些充分成分原因，在图 5.3 中以叉号标出。
\end{techpoint}
\renewcommand{\TPname}{技术注记}

\newcommand{\sccCross}[1]{%
  \draw[red,very thick] (-0.5,-0.5)--(0.5,0.5);
  \draw[red,very thick] (-0.5,0.5)--(0.5,-0.5);}

\begin{figure}[htbp]
\centering
\begin{tikzpicture}[causal,scale=1]
  \begin{scope}
    \sccOne{U_0=1}
  \end{scope}
  \begin{scope}[xshift=4.4cm]
    \sccTwo{A=1}{U_1=1}
  \end{scope}
  \begin{scope}[xshift=8.8cm]
    \sccTwo{A=0}{U_2=1}\sccCross{}
  \end{scope}
  \begin{scope}[yshift=-2.2cm]
    \sccTwo{U_3=1}{E=1}
  \end{scope}
  \begin{scope}[xshift=4.4cm,yshift=-2.2cm]
    \sccTwo{U_4=1}{E=0}\sccCross{}
  \end{scope}
  \begin{scope}[xshift=8.8cm,yshift=-2.2cm]
    \sccThree{A=1}{E=1}{U_5=1}
  \end{scope}
  \begin{scope}[yshift=-4.4cm]
    \sccThree{A=1}{E=0}{U_6=1}\sccCross{}
  \end{scope}
  \begin{scope}[xshift=4.4cm,yshift=-4.4cm]
    \sccThree{A=0}{E=1}{U_7=1}\sccCross{}
  \end{scope}
  \begin{scope}[xshift=8.8cm,yshift=-4.4cm]
    \sccThree{A=0}{E=0}{U_8=1}\sccCross{}
  \end{scope}
\end{tikzpicture}
\caption{效应单调时不可能存在的充分成分原因（以叉号标出）}
\label{fig:5.3}
\end{figure}
